\documentclass[11pt]{article}
\usepackage[margin=1in]{geometry}
\usepackage{booktabs}
\usepackage{amsmath}
\usepackage{graphicx}
\usepackage{hyperref}
\usepackage{xcolor}

\newcommand{\todo}[1]{\textcolor{red}{[TODO: #1]}}
\newcommand{\ours}{\textsc{Open-Jev}}

\title{Open System-One Decision Models:\\
Frozen Diffusion-LM Backbones and Specialist Decision Heads\\
for Calibrated Scientific Decisions}
\author{Anonymous}
\date{\today}

\begin{document}
\maketitle

\begin{abstract}
Structured decisions---selecting among options, verifying whether evidence
supports a claim, and grading response quality on ordinal scales---are the
atomic operations of automated scientific pipelines. Proprietary
``System-One'' models deliver such decisions as typed API calls without
generation, but are closed. We present an open, trainable counterpart: a
frozen masked-diffusion language model backbone (LLaDA-8B) with small
specialist decision heads producing probabilities for three decision types
(choice, noul, score). We contribute three findings. (1)~A deterministic
canonical ordering of options yields exact permutation invariance
(option-order flip rate $0.77 \to 0.00$), where train-time logit averaging
does not. (2)~Multi-layer attention pooling over option spans outperforms
last-layer mean pooling for choice ($+7$ points), while an ordinal auxiliary
loss improves graded scoring ($+9$ points). (3)~Scaling training data
$30\times$ on an out-of-domain corpus improves in-distribution accuracy to
$0.96$ but does not transfer ($-13$ points on the target domain), showing
that coverage of the target reasoning distribution---not volume---is the
binding constraint. On public benchmarks, \ours{} reaches \todo{bench\_eval
numbers} accuracy on GPQA and \todo{} on SciFact verification, within reach
of published System-One results on general classification.
\end{abstract}

\section{Introduction}
% Motivation: scientific pipelines need typed decisions, not generated text.
% Proprietary System-One models (Jev, TypeSafe AI) are closed; the underlying
% architecture is unclear. We build an open equivalent and study what matters.
\todo{intro: typed decisions vs generation; routing to System Two;
contributions list}

\section{Setup}
\paragraph{Decision types.}
We consider three typed decisions over a state (context) and a question:
\textbf{choice} selects one of $K$ answer options; \textbf{noul} decides
whether evidence supports a claim ($K{=}2$); \textbf{score} assigns an
ordinal grade $0 < 1 < 2$ (unrelated / related-but-wrong / correct).

\paragraph{Model.}
Each example is encoded as a single masked sequence: context, question, and
options separated by marker tokens. A frozen masked-diffusion backbone
(LLaDA-8B-Instruct) produces token representations; a specialist head per
decision type maps pooled option-span features to class logits. Two head
architectures are compared: a mean-pooled MLP and \emph{AttnPool}, a learned
query attending over each option span on concatenated hidden states from
layers $\{8,16,24,32\}$. Only the heads are trained ($\sim$2M parameters);
the backbone remains frozen.

\paragraph{Canonical ordering.}
Options are sorted by token-id content before encoding, and gold indices
remapped. The presentation order is thus a pure function of the option
\emph{set}: any input permutation collapses to the same sequence and
permutation invariance holds exactly, at training, calibration, and
evaluation.

\paragraph{Ordinal auxiliary loss.}
For score we add a CORAL-style term
$\sum_{j=1}^{K-1} \mathrm{BCE}\big(P(y \ge j),\,\mathbf{1}[g \ge j]\big)$
on cumulative softmax mass, respecting the grade order without changing the
$K$-way decode.

\paragraph{Metrics.}
Accuracy, expected calibration error (ECE), option-order flip rate (accuracy
change under input permutation), and \emph{automation@5\%err}: the fraction
of decisions that can be auto-accepted while keeping realized error $\le
5\%$---the operating metric for a decision layer.

\section{Experiments}
\subsection{Scientific decision battery}
% elite battery: passage-grounded scientific QA, typed by reasoning kind
% (factual, causal, multihop, negation, numerical, definitional);
% pid-disjoint splits; 2240/6720/6720 choice/noul/score train rows.
\todo{describe elite battery + v2 corpus (qa\_raw, 25K pids, 3 corpora)}

\subsection{Public benchmarks}
GPQA main (448 graduate-level questions) and diamond subset (198) as
choice; SciFact dev (340 claim--abstract pairs) as noul and 3-way score;
plus general classification (SST-2, AG News, Enron spam, Banking77) to
probe out-of-domain transfer. Benchmarks are eval-only; temperature is
fitted on the battery dev split (no benchmark leakage).

\section{Results}
\subsection{Ablations on the scientific battery}
\begin{table}[h]
\centering\small
\begin{tabular}{lcccc}
\toprule
 & choice & noul & score & flip (choice) \\
\midrule
Baseline (mixed MLP)        & 0.775 & 0.707 & 0.522 & 0.76 \\
$+$ orders$=4$              & 0.721 & 0.782 & 0.586 & --   \\
$+$ AttnPool                & 0.847 & 0.636 & 0.588 & 0.75 \\
$+$ both                    & 0.749 & 0.780 & 0.576 & 0.70 \\
\midrule
Per-type heads              & 0.876 & 0.710 & 0.574 & 0.77 \\
Per-type $+$ canon. $+$ ord. & \textbf{0.870} & \textbf{0.783} & \textbf{0.608} & \textbf{0.00} \\
\bottomrule
\end{tabular}
\caption{Ablation on the scientific battery (frozen LLaDA-8B). Canonical
ordering eliminates order sensitivity exactly; per-type specialist heads
plus AttnPool give the best choice accuracy and automation.}
\end{table}

Key findings: (1)~train-time logit averaging over permutations does
\emph{not} fix order sensitivity and costs choice accuracy;
(2)~canonical ordering achieves flip rate exactly $0$ at no accuracy cost;
(3)~AttnPool on intermediate layers helps choice ($+7$) but hurts noul,
motivating per-type heads; (4)~the ordinal auxiliary loss lifts score
$0.522 \to 0.608$.

\subsection{Data scaling does not transfer across domains}
\begin{table}[h]
\centering\small
\begin{tabular}{lcc}
\toprule
 & in-distribution (v2) & elite (target) \\
\midrule
choice & 0.965 & 0.738 \\
noul   & 0.942 & 0.782 \\
score  & --    & 0.616 \\
\bottomrule
\end{tabular}
\caption{Heads retrained on $30\times$ more data from a different QA
distribution (qa\_raw corpora) reach near-ceiling in-distribution accuracy
but do not transfer to the elite battery, whose questions stress
causal/multihop/negation reasoning absent from the source corpus.}
\end{table}

\subsection{Public benchmarks}
\begin{table}[h]
\centering\small
\begin{tabular}{lccc}
\toprule
benchmark & type & \ours{} & reference \\
\midrule
GPQA main       & choice & \todo{} & GPT-4 $\sim$0.39; expert $\sim$0.65 \\
GPQA diamond    & choice & \todo{} & -- \\
SciFact         & noul   & \todo{} & -- \\
SciFact         & score  & \todo{} & -- \\
\midrule
Enron spam      & choice & \todo{} & published System-One: 0.987 \\
SST-2           & choice & \todo{} & 0.957 \\
AG News         & choice & \todo{} & 0.913 \\
Banking77       & choice & \todo{} & 0.760 \\
\bottomrule
\end{tabular}
\caption{Zero-shot transfer to public benchmarks. Published System-One
numbers are third-party measurements (Kumar, 2026) on identical public
datasets; our scientific decision heads were trained on an unrelated
domain.}
\end{table}

\section{Discussion}
\todo{what a System-One layer buys: determinism, calibration, automation
budget; routing uncertain cases (auto@5\%); the domain-coverage lesson;
relation to published System-One results (third-party reports only).}

\section{Limitations}
Backbone pretraining may include benchmark-adjacent text (standard public
benchmark caveat). Our model is trained on one scientific QA distribution;
cross-domain transfer is limited by construction. Score remains the weakest
decision type for all systems studied, including the proprietary reference.

\bibliographystyle{plain}
\begin{thebibliography}{9}
\bibitem{llada} Nie et al. Large Language Diffusion Models. 2025.
\bibitem{gpqa} Rein et al. GPQA: A Graduate-Level Google-Proof Q\&A
Benchmark. COLM 2024.
\bibitem{scifact} Wadden et al. Fact or Fiction: Verifying Scientific
Claims. EMNLP 2020.
\bibitem{kumar} Kumar, A. Jev, measured. Third-party benchmark report, 2026.
\bibitem{corallike} Cao et al. Rank consistent ordinal regression for
neural networks. 2020.
\end{thebibliography}

\end{document}
